\section{Infinity, the Absolute, the Self}

There are two approaches to understand logical systems:

\begin{itemize}
\item \textbf{Axiomatic}: A formal symbolic system with definitions and axioms. The semantic meanings of the symbols are not necessary to understand the system; theorems are proved just by the manipulation of symbols. 
\item \textbf{Naïve}: The semantic meaning of the symbols are understood in their everyday meaning. 
\end{itemize}
The Axiomatic method is used by mathematicians and analytic philosophers. The Naïve method is used by metaphysicians and Hermetists, for who the meaning (semantics) is more important. Hence they can use the principle of analogy to derive higher principles than what the axioms alone can prove.  There are three levels of logic, each dealing with a different subject matter.

\begin{enumerate}
\item \textbf{Predicate of Formal Logic}: This deals with logical relationships, inferences, properties, numbers, sets, infinity, order, and so on. As such, it is confined to mental concepts, not to physical or existing things. For the metaphysician, its goal is an understanding of infinity. 
\item \textbf{Modal or Organic Logic}: This deals with necessity, possibility, contingency, which means it needs to account for existing things, not merely ideas and concepts. Its foal is knowledge of the Totality. 
\item \textbf{Deontic or Moral Logic}: This deals with obligation, permissibility, prohibition, commands, promises, and so on. As such, it needs to account for action in the world. Its goal is the I or Self. 
\end{enumerate}
Logic begins in the domain of intellectuality, but may be alchemically transformed into spirituality. Valentin Tomberg, in the Letter on the Fool, describes this process:

\begin{quotex}
in the domain of the individual's inner life, it is analogous to the work of spiritual alchemy which operates on the historical plane. This means to say that the individual soul begins initially with the experience of the separation and opposition of the spiritual and intellectual elements within it, then advances to—or resigns itself to—parallelism, i.e., a kind of ``peaceful coexistence" of these two elements within it. Subsequently it arrives at cooperation between spirituality and intellectuality which, proving to be fruitful, eventually becomes the complete fusion of these two elements in a third element: the ``philosopher's stone" of the spiritual alchemy of Hermeticism. The beginning of this final stage is announced by the fact that logic becomes transformed from formal logic, passing through the intermediary stage of organic logic, into moral logic. 

\end{quotex}
We dealt with Predicate logic in \textit{Naïve Set Theory}\footnote{See Section \ref{sec:NaiveSetTheory} in this book.}. Now we deal with Organic and Moral logic. This will show that the world process can be understood on four levels:

\begin{enumerate}
\item What is a possibility of manifestation 
\item What is necessary in the world 
\item What is commanded by God 
\item What is willed by God 
\end{enumerate}
\paragraph{Modal or Organic Logic}
Modal logic deals with the following modes of being:

\begin{itemize}
\item Necessary: What is necessary 
\item Possible: What is possible 
\item Impossible: what cannot be manifested 
\item Contingent: what depends on something else (compossible) 
\end{itemize}
As such, these all deal with the possibilities of manifestation, i.e., the Infinity of God as the beyond all the infinities of formal logic. The ``impossible" are possibilities of non-manifestation. Contingent things are possibilities that can manifest only under specific conditions. In philosophical terms, necessary things must exist in all possible worlds. Better said from the metaphysical point of view, it exists in all degrees of existence.

Modal logic discovers necessary connections among the things of the world. A system can be built up through ever increasing levels of abstraction, so that more is known from fewer and fewer principles. Ultimately, this leads to an understanding of the Absolute. In these systems of Absolute Idealism, the whole is greater than its parts. For example, the community is greater than any individual.

Deontic logic also comprises mereology, which is the study of how parts are related to the whole and systems theory, which studies how systems contribute to the whole.

\paragraph{Deontic or Moral Logic}
Deontic logic deals with moral agents acting in the world. Their acts fall into four categories of what is:

\begin{itemize}
\item \textbf{Obligatory}: actions that are commanded or obligatory 
\item \textbf{Permissible}: actions that are permitted 
\item \textbf{Prohibited}: actions that are prohibited 
\item \textbf{Facultative}: actions that are morally indifferent 
\end{itemize}
Tomberg illustrates how an axiom transforms through the degrees of logic:

\begin{itemize}
\item For formal logic, it is an axiom the whole is greater than its parts. That is because it deals with quantities, so it is true within that domain. 
\item Organic logic deals with functions, not quantities. Hence, although the heart is quantitatively a small part of the body, its function is necessary for the body to exist as a whole. This requires a new axiom: the part can be equal to the whole. 
\item Moral logic deals with values, hence the axiom is transformed to: the part can be greater than the whole. Hence, it is immoral to sacrifice one man for the good of the whole. 
\end{itemize}
Morality cannot be reduced to a mechanical process of reasoning, which is the goal of deontic logic. For it to be transformed into moral logic, one must live a life of prayer and meditation. One discovers the axioms that make moral logic even possible. Tomberg points to Rene Descartes and Immanuel Kant as discoverers of those axioms.

\subparagraph{Descartes}
Tomberg looks past the logical, philosophical, and psychological criticisms of Descartes in order to get at Descartes' actual experience, which is beyond such critiques. That is the awareness of the Transcendent Self that transcends both Thought and Extension. In other words, he came to the self-realization as the ``thinker of thoughts", leading to complete certainty of the ``I AM". The axiom, then, is:

\begin{itemize}
\item The Reality of the Self 
\end{itemize}
\subparagraph{Kant}
Kant elucidated the following three additional axioms, without which, there is no possibility of moral action:

\begin{itemize}
\item The Reality of God or infinite perfectibility 
\item The Reality of Moral Freedom or morality as such 
\item The Reality of the immortality of the soul or the possibility of infinite perfectioning
\end{itemize}


\flrightit{Posted on 2021-06-03 by Cologero }
